{\scriptsize
    \begin{minipage}[t]{0.49\textwidth}
        \begin{tabularx}{\linewidth}{|p{0.35\textwidth}|X|}
            \hline 
            \textbf{Notions et contenus} & \textbf{Capacités exigibles} \\
            \hline 
            \multicolumn{2}{|l|}{4.3. Mécanique des fluides} \\
            \hline 
            \multicolumn{2}{|l|}{4.3.1. Description d'un fluide en mouvement} \\
            \hline 
             Champ eulérien des vitesses. Lignes de champ. Tubes de champ. & 
              Définir et utiliser l'approche eulérienne.  \\
            \hline 
            Écoulement stationnaire.  & Discuter du caractère stationnaire d'un écoulement en fonction du référentiel d'étude.\\
            \hline 
            Dérivée particulaire de la masse volumique. Écoulement incompressible. & Établir l'expression de la dérivée particulaire de la masse volumique. Utiliser l'expression de la dérivée particulaire de la masse volumique pour caractériser un écoulement incompressible. \\
            \hline
             Débit massique. Débit volumique.  &  Définir le débit massique et l'écrire comme le flux du vecteur densité de courant de masse à travers une surface orientée. Définir le débit volumique et l'écrire comme le flux du champ de vitesse à travers une surface orientée. \\
            \hline
        \end{tabularx}
\end{minipage}
\begin{minipage}[t]{0.49\textwidth}
    \begin{tabularx}{\linewidth}{|p{0.35\textwidth}|X|}
        \hline 
        Équation locale de conservation de la masse. & Établir l'équation locale de conservation de la masse dans le seul cas d'un problème unidimensionnel en géométrie cartésienne. Citer et utiliser une généralisation admise en géométrie quelconque à l'aide de l'opérateur divergence et son expression fournie. \\
        \hline 
        Caractérisation d'un écoulement incompressible par la divergence du champ des vitesses. & Traduire localement, en fonction du champ de vitesses, le caractère incompressible d'un écoulement. \\
        \hline 
        Dérivée particulaire du champ de vitesse : terme local ; terme convectif. & Associer la dérivée particulaire de la vitesse à l'accélération de la particule de fluide qui passe en un point.
        Utiliser l'expression de l'accélération, le terme convectif étant écrit sous la forme (v.grad) $\mathbf{v .}$
        
        Utiliser l'expression fournie de l'accélération convective en fonction de grad ( $v^2 / 2$ ) et rot $\mathbf{v} \times \mathrm{v}$. \\
        \hline
        Écoulement irrotationnel défini par la nullité du rotationnel du champ des vitesses en tout point ; potentiel des vitesses. & Traduire localement, en fonction du champ de vitesses, le caractère irrotationnel d'un écoulement et en déduire l'existence d'un potentiel des vitesses. \\
        \hline
    \end{tabularx}
\end{minipage}
}
